% 07_conclusion.tex
\section{Conclusion}
\label{sec:conclusion}

We have reproduced the Q4 pLDDT-fragmentation analysis of Cazals \&
Sarti~\cite{cazals2025alphafold} at the scale of the \textit{H.~sapiens}
proteome using AlphaFold-DB v4, implementing the entire pipeline from
scratch in Python with no dependence on the authors' code.
The path-graph filtration, Union-Find with the Elder Rule, all five
topological statistics, and the PLM computation via GUDHI are each
validated against a null random model and three prototype proteins before
being applied proteome-wide.

The central reproduction target --- a Pearson correlation of $r \approx
0.97$ between $f^+_{cp}$ and $H_p$ --- is recovered at $r = 0.928$
across 16{,}787 F1 fragments with $n \geq 200$ residues.
The residual gap of $0.042$ is systematically traced to the
floating-point pLDDT precision of AlphaFold-DB v4, which produces
$\sim\!10\times$ more distinct filtration levels per protein than the
effectively integer-valued distributions used in the paper.
The same mechanism halves the number of fragmentation candidates
identified by the three-way filter ($n \geq 200$, $H_p \geq 0.25$,
PLM $\geq 3$ at $t_p = 0.025$): 42 proteins vs.\ the paper's
$\approx 86$.
We quantify both gaps and explain them from first principles, converting
what could be seen as a failure to reproduce into a precise, falsifiable
hypothesis about the evolution of AlphaFold-DB data precision.

Beyond reproduction, the novel Q1$\times$Q4 cross-correlation analysis
reveals a 3.08-fold enrichment of fragmented proteins in the arity bin
$(a_{25}=7, a_{75}=13)$ (Fisher exact test, $p = 1.76 \times 10^{-7}$).
This finding --- derived from 3D $\mathrm{C}_\alpha$ packing geometry
independently of any pLDDT information --- provides orthogonal structural
evidence that the PLM $\geq 3$ filter identifies a genuine biological
phenotype: multi-domain proteins with compact individual domains separated
by flexible linkers.
The convergence of two independent descriptors on the same protein set
strengthens the biological credibility of the topological fragmentation
signal.
